\documentclass[11pt]{article}

% --------------------
% Packages
% --------------------
\usepackage[margin=1in]{geometry}
\usepackage{hyperref}
\usepackage{enumitem}
\usepackage{array}
\usepackage{booktabs}
\usepackage{xcolor}
\usepackage{listings}
\usepackage{courier}
\usepackage{titlesec}
\usepackage{setspace}

% --------------------
% Global spacing tweaks
% --------------------
\setstretch{1.15}                 % overall line spacing
\setlist[itemize]{itemsep=0.4em}  % bullet spacing

% --------------------
% Section formatting
% --------------------
\titleformat{\section}
  {\large\bfseries}
  {\thesection.}
  {0.75em}
  {}

\titleformat{\subsection}
  {\normalsize\bfseries}
  {}
  {0em}
  {}

\titlespacing*{\section}{0pt}{2.0em}{1.0em}
\titlespacing*{\subsection}{0pt}{1.2em}{0.6em}

% --------------------
% Code formatting
% --------------------
\lstset{
  basicstyle=\ttfamily\small,
  columns=fullflexible,
  frame=single,
  breaklines=true,
  showstringspaces=false,
  xleftmargin=1em,
  xrightmargin=1em
}

% --------------------
% Document
% --------------------
\begin{document}

\begin{center}
{\Large \textbf{CS 147 — Verilog Usage Rules}}\\[0.75em]
Spring 2026\\
\vspace{0.5em}
\textit{Adapted from CS/ECE 552 Verilog Rules}\\
\textit{Original author: Andy Phelps (CS 552, Spring 2006)}\\
Last Modified: January 2026
\end{center}

\vspace{1.5em}

\tableofcontents
\newpage

% ==========================================================
\section{Introduction}

In \textbf{CS 147}, you are restricted to a carefully chosen subset of the
Verilog hardware description language. This document defines and explains
those restrictions.

This is \emph{not} a complete Verilog reference. Instead, it is intended
to guide you toward writing Verilog that closely reflects real hardware
structures and design practices.

For full language references, consult:
\begin{itemize}
  \item Online Verilog HDL Quick Reference Guide
  \item ASIC World Verilog Quick Reference
\end{itemize}

% ==========================================================
\section{Justification for These Limitations}

The goal of these restrictions is to force you to \textbf{think like a
hardware designer}, not a software programmer.

While Verilog can be used in a highly abstract, C-like manner, doing so
provides little educational value in an introductory architecture course.

Limiting the language:
\begin{itemize}
  \item Encourages schematic-accurate thinking
  \item Prevents simulation-order–dependent bugs
  \item Produces designs that synthesize reliably
  \item Mirrors constraints commonly imposed in industry
\end{itemize}

Design your Verilog as if you were verbally describing real hardware
components and their connections.

% ==========================================================
\section{Allowed Keywords}

Verilog keywords for this course fall into three categories:
\begin{itemize}
  \item \textbf{Always Allowed}
  \item \textbf{Allowed with Stipulations}
  \item \textbf{Never Allowed}
\end{itemize}

This list is not guaranteed to be exhaustive. When in doubt, ask a TA
\emph{before} using a construct.

\subsection*{Keywords Allowed with Stipulations}

The following rules apply \textbf{only} to keywords in this category:

\begin{itemize}
  \item \texttt{case}, \texttt{casex}
    \begin{itemize}
      \item A \texttt{default} clause is required
      \item If the default executes, an error must be asserted
      \item Large case statements (\(\sim\)12 lines or more) must be placed in their own module
    \end{itemize}

  \item \texttt{reg}
    \begin{itemize}
      \item May only be used for outputs of a case statement
    \end{itemize}

  \item \texttt{always}
    \begin{itemize}
      \item May only be used to introduce a case statement
    \end{itemize}

  \item \texttt{begin} / \texttt{end}
    \begin{itemize}
      \item May only wrap the body of a case statement
    \end{itemize}
\end{itemize}

% ==========================================================
\section{Allowed Operators}

You are restricted to simple boolean and structural operators.

\begin{itemize}
  \item Shift operators are allowed \textbf{only} with constant shift amounts
  \item The ternary operator (\texttt{a ? b : c}) is allowed and encouraged
  \item Concatenation (\texttt{\{a,b\}}) and reduction operators are allowed
  \item Equality (\texttt{==}) and inequality (\texttt{!=}) are permitted
\end{itemize}

\subsection*{Explicitly Forbidden}

The following categories of operators are \textbf{not allowed}:
\begin{itemize}
  \item Arithmetic: \texttt{+ - * / \%}
  \item Relational: \texttt{< > <= >=}
\end{itemize}

\subsection*{Operator Summary}

\begin{center}
\renewcommand{\arraystretch}{1.25}
\begin{tabular}{>{\ttfamily}l l | >{\ttfamily}l l}
\toprule
\multicolumn{2}{c|}{\textbf{Allowed}} & \multicolumn{2}{c}{\textbf{Not Allowed}} \\
\midrule
\textasciitilde m & Inversion      & m + n & Addition \\
m \& n            & Bitwise AND    & m - n & Subtraction \\
m | n             & Bitwise OR     & m * n & Multiplication \\
m \textasciicircum n & Bitwise XOR & m / n & Division \\
\&m               & AND Reduction  & m \% n & Modulus \\
m == n            & Equality       & m \&\& n & Logical AND \\
m != n            & Inequality     & m || n & Logical OR \\
m << const        & Shift Left     & m < n & Less Than \\
m >> const        & Shift Right    & m > n & Greater Than \\
test ? m:n        & Ternary        &       & \\
\bottomrule
\end{tabular}
\end{center}

% ==========================================================
\section{Usage Examples}

All hardware designs in this course consist of:
\begin{itemize}
  \item \textbf{Combinational logic}
  \item \textbf{State}, implemented using D flip-flops
\end{itemize}

All state must be held in provided D flip-flop modules with synchronous
reset. Do \emph{not} create custom flip-flops using procedural tricks.

\subsection*{Combinational Logic}

Combinational logic should primarily be expressed using \texttt{assign}
statements:

\begin{lstlisting}
assign xyz = a & b;
\end{lstlisting}

\subsection*{Module Structure}

Modules should follow a clear, consistent structure:

\begin{lstlisting}
module my_module_name(param1, param2);
  input  param1;
  output param2;

  wire [7:0] signalName;

  assign signalName = expression;

  submodule instance(
    .input0(param1),
    .output0(param2)
  );
endmodule
\end{lstlisting}

\subsection*{Case Statement Example}

\begin{lstlisting}
always @* case (s)
  2'b00: out = i0;
  2'b01: out = i1;
  2'b10: out = i2;
  2'b11: out = i3;
endcase
\end{lstlisting}

\textbf{Important:} Every output must be assigned in every case. Failure
to do so will infer latches and produce incorrect hardware.

% ==========================================================
\section{Verilog Check Program}

Your designs will be scanned using the \texttt{Vcheck} program to detect
illegal constructs. You are required to run it and submit the output.

Available methods:
\begin{itemize}
  \item \texttt{java Vcheck < myfile.v}
  \item \texttt{vcheck.sh < myfile.v}
  \item \texttt{vcheck-all.sh}
\end{itemize}

\end{document}
